\documentclass[a4paper]{article}

\title{C/C++-Kurs für Kinder ab der 5. Klasse}
\author{Jürgen ``George'' Sawinski}

\usepackage{polyglossia}
\setdefaultlanguage{german}
\setotherlanguage{english}\usepackage{graphicx} % allow embedded images
  \setkeys{Gin}{width=\linewidth,totalheight=\textheight,keepaspectratio}
  \graphicspath{{images/}} % set of paths to search for images

\usepackage{newtxtext}
\usepackage{newtxmath}
\usepackage{couriers}

\usepackage{amsmath}
\usepackage{units}    % non-stacked fractions and better unit spacing
\usepackage{multicol} % multiple column layout facilities
\usepackage{fancyvrb} % extended verbatim environments
  \fvset{fontsize=\normalsize}% default font size for fancy-verbatim environments
\usepackage{hyperref}
%   \hypersetup{
%     pdfborder=1 1 1,
%     colorlinks=false}
\usepackage{tikz}
  \usetikzlibrary{positioning, fit, backgrounds}

\usepackage{c-kurs}

\begin{document}
\maketitle

    \begin{abstract}\noindent
    Dieser Kurs ist gedacht für Kinder ab der Grundschule mit guten mathematischen Fähigkeiten und dem Drang, programmieren zu lernen. Er entstand, nachdem die Söhne des Autors ihn über ein Jahr beknieten, ein Spiel zu entwickeln.

    Der erste Teil des Kurses beschäfigt sich damit, alle grundlegenden Kenntnisse zu erwerben, um eine Konsolen-basierte Regensimulation zu in C zu programmieren, was auch mathematische und physikalische Grundlagen beinhaltet.

    Der zweite Teil des Kurses wird das Ziel haben, das Spiel Tetris® in C++ umzusetzen, wobei das Hauptaugenmerk dabei auf die Grundlagen Objekt-orientierter Programmierung gelegt wird.
    \end{abstract}

\tableofcontents

\section*{Vorwort}

Dieser Kurs ist aus der gemeinsamen Arbeit des Autors mit seinen Söhnen entstanden.
Die Idee, einen C/C++-Kurs für Kinder zu schreiben, kam von den Jungs selbst,
die unbedingt programmieren lernen wollten.

Der Autor hat die Inhalte erarbeitet, erklärt und in LaTeX festgehalten.
Bei der Strukturierung des Materials und der Erstellung einiger LaTeX-Makros
wurde er durch \textbf{DeepSeek}, ein KI-gestütztes Werkzeug, unterstützt.

DeepSeek lieferte Vorschläge zur Gliederung, zur didaktischen Aufbereitung
und zu technischen Umsetzungen in LaTeX. Alle Vorschläge wurden vom Autor
geprüft, angepasst und auf die Bedürfnisse seiner Kinder zugeschnitten.

\textit{-- Jürgen ``George'' Sawinski, \today}

\newpage

%%%%%%%%%%%%%%%%%%%%%%%%%%%%%%%%%%%%%%%%%%%%%%%%%%%%%%%%%%%%%%%%%%%%%%%%%%%%%%%%%%%
%%%%%%%%%%%%%%%%%%%%%%%%%%%%%%%%%%%%%%%%%%%%%%%%%%%%%%%%%%%%%%%%%%%%%%%%%%%%%%%%%%%
%%%%%%%%%%%%%%%%%%%%%%%%%%%%%%%%%%%%%%%%%%%%%%%%%%%%%%%%%%%%%%%%%%%%%%%%%%%%%%%%%%%
\section{Etwas Mathematik}
\label{sec:etwas-mathematik}

Computer rechnen im sogenannten Binärsystem, sie kennen nur 0 und 1. Um größere Zahlen darstellen zu können, werden mehrere dieser Nullen und Einsen, genannt ``\textbf{Bit}'' aneinandergereit. Am gebräuchlichsten dabei ist das ``\textbf{Byte}'' (siehe Abb.~\ref{fig:byte_bits}), das aus 8 Bit besteht.

\begin{figure}[h]
  \centering
  \begin{tikzpicture}[
      cell/.style={
          draw,
          rectangle,
          minimum width=0.8cm,
          minimum height=0.8cm,
          font=\ttfamily\large,
          outer sep=0pt,          % Kein Abstand zwischen Zellen
          anchor=base,
      }
  ]

  % Bit-Positionen (0 bis 7) nebeneinander
  \foreach \i in {0,1,...,7} {
      \node[cell] at (7-\i,0) (bit\i) {\i};
  }

  % Border around entire byte
  \node[draw=black, thick, inner sep=4pt, fit=(bit0)(bit7)] (bytebox) {};

  % Etikett "Bit" darüber
  \node[font=\sffamily, left=0.1cm of bytebox] {\textbf{Bit}};

  \end{tikzpicture}
  \caption{Ein Byte mit den Bit-Positionen 0 bis 7.}
  \label{fig:byte_bits}
\end{figure}

\noindent
Jedem Bit wird nun ein Wert zugeordnet. Das erfolgt mathematisch gemäß:
\begin{equation*}
  2^b
\end{equation*}
wobei $b$ die Nummer des Bits ist. Das ist eine sogenannte ``Potenz'', eine mathematische Operation die sagt, wie oft eine Zahl mit sich selbst multipliziert wird. In diesem Fall wird die ``2'' $b$-mal mit sich selbst multipliziert, zum Beispiel:
\begin{align*}
  2^1 &= 2              &= 2\\
  2^2 &= 2\cdot2        &= 4\\
  2^3 &= 2\cdot2\cdot2  &= 8
\end{align*}
Interessant ist nun die Frage, was $2^0$ ist. Wenn ihr genau hinschaut, dann merkt ihr, dass, um in die nächste Zeile zu kommen, ihr das Ergebnis mit $2$ multipliziern müsst. Umgekehrt genauso: um in die vorhergehende Zeile zu kommen, muss man durch $2$ teilen. Wenn also $2^1$ zwei ist und man, um zu $2^0$ zu kommen durch zwei Teilen muss, dann ergibt sich
\begin{equation*}
  2^0 = 1\,.
\end{equation*}
Übrigens, jede Zahl hoch $0$ hat als Ergebnis die $1$ (wir kommen später darauf zurück).

Um auf das \textbf{Byte} zurückzukommen, heißt dass, dass den \textbf{Bits} jeweils die Werte 1 bis 128 zugeordnet sind (siehe Abb.~\ref{fig:byte_values}).

\begin{figure}[h]
  \centering
  \begin{tikzpicture}[
      cell/.style={
          draw,
          rectangle,
          minimum width=0.8cm,
          minimum height=0.8cm,
          font=\ttfamily\large,
          outer sep=0pt,          % Kein Abstand zwischen Zellen
          anchor=base,
      }
  ]

  % Zeile 1: Bit-Positionen (0 bis 7)
  \foreach \i in {0,1,...,7} {
      \node[cell] at (7-\i,0) (bit\i) {\i};
  }

  % Zeile 2: Werte (1, 2, 4, 8, 16, 32, 64, 128)
  \foreach \i in {0,1,...,7} {
      \node[cell] at (7-\i,-1.2) (val\i) {\pgfmathparse{int(2^\i)}\pgfmathresult};
  }

  % Border around
  \node[draw=black, thick, inner sep=4pt, fit=(bit0)(bit7)] (bytebox) {};
  \node[draw=black, thick, inner sep=4pt, fit=(val0)(val7)] (valbox) {};


  % Etiketten "Bit" und "Wert" darüber
  \node[font=\sffamily, left=0.1cm of bytebox] {\textbf{Bit}};
  \node[font=\sffamily, left=0.1cm of valbox] {\textbf{Wert}};

  \end{tikzpicture}
  \caption{Ein Byte mit Bit-Positionen und den zugehörigen Werten ($2^b$).}
  \label{fig:byte_values}
\end{figure}

\noindent
Damit können wir nun die Zahlen 0 bis 255 darstellen, z.B.:
\begin{align*}
  00000000 &\rightarrow 0\,, \\
  00000011 &\rightarrow 3\,, \\
  01011000 &\rightarrow 88\,, \\
  11111111 &\rightarrow 255\,.
\end{align*}
Das Addieren funktioniert nun genauso, wie ihr es aus der Schule, in der natürlich das Zehnersystem verwendet wird, gelernt habt. Beispiel:
\begin{lstlisting}
      00100111       39
    + 00000110     +  6
          11         1     (Übertrag)
     ---------      ---
      00101101       45
\end{lstlisting}
Entsprechend lässt sich auch die Multiplikation durchführen:
\begin{lstlisting}
      00001010·00101010 =     10·42

      00000000
       00000000
        00000000
         00000000
          00101010
           00000000
            00101010
             00000000
             1 1             (Übertrag)
      ---------------
            110100100        420
\end{lstlisting}
Wenn ihr genau hinseht, dann bemerkt ihr, dass wir zwei Bytes multipliziert haben, aber das Ergebnis benötigt ein Bit mehr. Das führt zu dem Problem, dass, wenn ihr auf einem Computer diese Multiplikation durchführt, das höchste Bit (die letzte ``1'' links) abgeschnitten wird und das Ergebnis, dass der Computer euch liefert, $164$ ist. Das gleich passiert bei einem Byte, wenn ihr zu $255$ die $1$ dazuzählt, das Ergebnis ist $0$!

Computer können auch negative Zahlen darstellen. Dazu nützt man aus, dass, wenn $255 + 1$ zum Ergebnis $0$ hat, dann hat $0-1$ als Ergebnis $-1$, binär also 11111111. Dies nennt man das ``Zweierkomplement'' (den Begriff müsst ihr euch natürlich nicht merken). Wichtig ist: dass höchste Bit zeigt, dass es sich um eine negative Zahl handelt. Damit können dann mit einem Byte die Zahlen -128 bis 127 dargestellt werden:
\begin{lstlisting}
    10000000        -128
    10000001        -127
    ...
    11111111        -  1
    00000000           0
    00000001           1
    ...
    01111111         127
\end{lstlisting}
Wie wird nun aus einer positiven Zahl eine Negative? Das passiert, indem man alle Bits invertiert (das heißt, aus 0 wird 1 und aus 1 eine 0) und 1 addiert:
\begin{lstlisting}
(-127)  10000001  ->  01111110
                     +00000001
                      --------
                      01111111     (127)
\end{lstlisting}

%%%%%%%%%%%%%%%%%%%%%%%%%%%%%%%%%%%%%%%%%%%%%%%%%%%%%%%%%%%%%%%%%%%%%%%%%%%%%%%%%%%
%%%%%%%%%%%%%%%%%%%%%%%%%%%%%%%%%%%%%%%%%%%%%%%%%%%%%%%%%%%%%%%%%%%%%%%%%%%%%%%%%%%
%%%%%%%%%%%%%%%%%%%%%%%%%%%%%%%%%%%%%%%%%%%%%%%%%%%%%%%%%%%%%%%%%%%%%%%%%%%%%%%%%%%
\section{Hallo Welt!}
\label{sec:hallo-welt}

In alter Tradition ist das erste Programm, das man zu sehen bekommt, wenn man eine neue Programmiersprache lernt, gewöhnlich das ``Hallo Welt!'' Programm \cite{wikipedia:hallo-welt}. Da wir uns zunächst auf die Programmiersprache ``C'' beschränken, ist das:

\begin{code}
/* Das ``Hallo Welt''-Programm. */
#include <stdio.h>

int main() {
  printf("Hallo Welt!\n");
  return 0;
}
\end{code}

\section*{Zeile für Zeile erklärt}

Schauen wir uns das Programm Schritt für Schritt an:

\begin{description}
    \item[\texttt{/* Das ``Hallo Welt''-Programm. */}]
        Das ist ein \textbf{Kommentar}.
        Der Computer ignoriert ihn völlig. Er ist nur für uns Menschen da,
        um den Code zu erklären.

        In C gibt es zwei Arten von Kommentaren:
        \begin{itemize}
            \item \texttt{/* ... */} — für längere Kommentare über mehrere Zeilen.
            \item \texttt{// ...} — für kurze Kommentare in einer Zeile.
        \end{itemize}

    \item[\texttt{\#include <stdio.h>}]
        Das ist eine \textbf{Header-Datei} (auf Deutsch auch \emph{Kopfdatei}).

        \texttt{stdio.h} steht für \textit{Standard Input Output}.
        Sie enthält Informationen über Funktionen wie \texttt{printf},
        die wir später benutzen wollen.

        \textbf{Merke:} Wenn wir eine Funktion benutzen wollen, die wir nicht
        selbst geschrieben haben, müssen wir die dazugehörige Header-Datei
        einbinden.

    \item[\texttt{int main() \{ }]
        Das ist die \textbf{Hauptfunktion} (englisch: \textit{main function}).

        Jedes C-Programm startet hier. Der Computer führt den Code in den
        geschweiften Klammern \texttt{\{ ... \}} aus.

        \texttt{int} bedeutet, dass die Funktion eine ganze Zahl (Integer)
        zurückgibt — das ist der Rückgabewert, der weiter unten kommt.

    \item[\texttt{printf("Hallo Welt!$\backslash$n");}]
        Das ist die \textbf{Ausgabe}.

        \texttt{printf} steht für \textit{print formatted} und schreibt Text
        auf den Bildschirm (genauer: in die Konsole).

        Das \texttt{$\backslash$n} ist ein \textbf{Zeilenumbruch}.
        Es sorgt dafür, dass der Cursor nach der Ausgabe in die nächste Zeile springt.

        \textbf{Wichtig:} Jede Anweisung in C endet mit einem Semikolon \texttt{;}.

    \item[\texttt{return 0;}]
        Das ist die \textbf{Rückgabe}.

        Sie sagt dem Betriebssystem: \textit{``Alles ist gut gelaufen.''}

        \texttt{0} bedeutet in der Regel \textit{``kein Fehler''}.
        Wenn ein Programm mit \texttt{return 0;} endet, weiß der Computer,
        dass das Programm erfolgreich war.

    \item[\texttt{\}} ]
        Die schließende geschweifte Klammer beendet die \texttt{main}-Funktion.
        Alles zwischen \texttt{\{} und \texttt{\}} gehört zur Funktion.
\end{description}

\section*{Zusammenfassung}

\begin{itemize}
    \item Ein C-Programm beginnt immer in der \texttt{main()}-Funktion.
    \item \texttt{printf()} schreibt Text auf den Bildschirm.
    \item \texttt{$\backslash$n} macht einen Zeilenumbruch.
    \item \texttt{return 0;} sagt dem Computer: ``Alles OK.''
    \item Jede Anweisung endet mit \texttt{;}.
\end{itemize}

\section*{Aufgaben}

\begin{enumerate}
    \item Ändere das Programm so, dass es \texttt{``Guten Tag!''} ausgibt.
    \item Ändere das Programm so, dass es deinen Namen ausgibt.
    \item Ändere das Programm so, dass es zwei Zeilen ausgibt:
        \begin{verbatim}
        Hallo!
        Welt!
        \end{verbatim}
    \item (Für Profis) Was passiert, wenn du das \texttt{$\backslash$n} weglässt?
\end{enumerate}


%%%%%%%%%%%%%%%%%%%%%%%%%%%%%%%%%%%%%%%%%%%%%%%%%%%%%%%%%%%%%%%%%%%%%%%%%%%%%%%%%%%
%%%%%%%%%%%%%%%%%%%%%%%%%%%%%%%%%%%%%%%%%%%%%%%%%%%%%%%%%%%%%%%%%%%%%%%%%%%%%%%%%%%
%%%%%%%%%%%%%%%%%%%%%%%%%%%%%%%%%%%%%%%%%%%%%%%%%%%%%%%%%%%%%%%%%%%%%%%%%%%%%%%%%%%
\newpage
\section{Zahlen}
\label{sec:zahlen}

\begin{code}
#include <stdio.h>

int main() {
  int zahl = 42;
  printf("Die Zahl: %d\n", zahl);
  return 0;
}
\end{code}

%%%%%%%%%%%%%%%%%%%%%%%%%%%%%%%%%%%%%%%%%%%%%%%%%%%%%%%%%%%%%%%%%%%%%%%%%%%%%%%%%%%
%%%%%%%%%%%%%%%%%%%%%%%%%%%%%%%%%%%%%%%%%%%%%%%%%%%%%%%%%%%%%%%%%%%%%%%%%%%%%%%%%%%
%%%%%%%%%%%%%%%%%%%%%%%%%%%%%%%%%%%%%%%%%%%%%%%%%%%%%%%%%%%%%%%%%%%%%%%%%%%%%%%%%%%
\newpage
\section{Fibonacci}
\label{sec:fibonacci}

\begin{code}
#include <stdio.h>

int main() {
  int a = 0;
  int b = 1;

  for(int i = 1; i <= 16; i = i+1) {
    int next = a + b;
    a = b;
    b = next;

    printf("Fibonacci %d: %d\n", i, a);
  }
  return 0;
}
\end{code}

%%%%%%%%%%%%%%%%%%%%%%%%%%%%%%%%%%%%%%%%%%%%%%%%%%%%%%%%%%%%%%%%%%%%%%%%%%%%%%%%%%%
%%%%%%%%%%%%%%%%%%%%%%%%%%%%%%%%%%%%%%%%%%%%%%%%%%%%%%%%%%%%%%%%%%%%%%%%%%%%%%%%%%%
%%%%%%%%%%%%%%%%%%%%%%%%%%%%%%%%%%%%%%%%%%%%%%%%%%%%%%%%%%%%%%%%%%%%%%%%%%%%%%%%%%%
\newpage
\section{Reihen}
\label{sec:reihen}

\begin{code}
#include <stdio.h>

int main() {

  return 0;
}
\end{code}

%%%%%%%%%%%%%%%%%%%%%%%%%%%%%%%%%%%%%%%%%%%%%%%%%%%%%%%%%%%%%%%%%%%%%%%%%%%%%%%%%%%
%%%%%%%%%%%%%%%%%%%%%%%%%%%%%%%%%%%%%%%%%%%%%%%%%%%%%%%%%%%%%%%%%%%%%%%%%%%%%%%%%%%
%%%%%%%%%%%%%%%%%%%%%%%%%%%%%%%%%%%%%%%%%%%%%%%%%%%%%%%%%%%%%%%%%%%%%%%%%%%%%%%%%%%
\appendix
\section{Software}

Als Programmierumgebung wurde \href{https://code.visualstudio.com/}{VSCode}

\section{Leitfaden}
\label{sec:leitfaden}

\subsection{C}
\label{sec:reference:c}

\subsubsection{Grundlagen}

%%%%%%%%%%%%%%%%%%%%%%%%%%%%%%%%%%%%%%%%%%%%%%%%%%%%%%%%%%%%%%%%%%%%%%%%%%%%%%%%%%%
\paragraph{Datentypen}

\subparagraph{Ganzzahlen}



\subparagraph{Fließkommazahlen}

%%%%%%%%%%%%%%%%%%%%%%%%%%%%%%%%%%%%%%%%%%%%%%%%%%%%%%%%%%%%%%%%%%%%%%%%%%%%%%%%%%%
\paragraph{Strukturen}

%%%%%%%%%%%%%%%%%%%%%%%%%%%%%%%%%%%%%%%%%%%%%%%%%%%%%%%%%%%%%%%%%%%%%%%%%%%%%%%%%%%
\paragraph{Schleifen}

%%%%%%%%%%%%%%%%%%%%%%%%%%%%%%%%%%%%%%%%%%%%%%%%%%%%%%%%%%%%%%%%%%%%%%%%%%%%%%%%%%%
\paragraph{Funktionen}

%%%%%%%%%%%%%%%%%%%%%%%%%%%%%%%%%%%%%%%%%%%%%%%%%%%%%%%%%%%%%%%%%%%%%%%%%%%%%%%%%%%
\paragraph{Schleifen}

%%%%%%%%%%%%%%%%%%%%%%%%%%%%%%%%%%%%%%%%%%%%%%%%%%%%%%%%%%%%%%%%%%%%%%%%%%%%%%%%%%%
\paragraph{Referenz}

\begin{man}[stdio]{printf}
    \begin{synopsis}
        \begin{plaincode}[gobble=8]
        #include <stdio.h>
        int printf(const char *format, ...);
        \end{plaincode}
    \end{synopsis}

    \begin{desc}
        \codeinline{printf()} schreibt formattierten Text in die Standardausgabe (stdout). Der \codeinline{format}-String enthält:
        \begin{itemize}
            \item Normale Zeichen (werden direkt ausgegeben)
            \item Format-Specifiers (beginnen mit \codeinline{\%})
        \end{itemize}

        \textbf{Die wichtigsten Format-Specifiers:}

        \begin{tabular}{ll}
        \texttt{\%d} & Ganzzahl (dezimal) \\
        \texttt{\%f} & Gleitkommazahl \\
        \texttt{\%c} & Zeichen (char) \\
        \texttt{\%s} & Zeichenkette (String) \\
        \texttt{\%x} & Hexadezimal (klein) \\
        \texttt{\%X} & Hexadezimal (groß) \\
        \texttt{\%p} & Zeiger (Pointer) \\
        \end{tabular}
    \end{desc}

    \textbf{Beispiel:}
    \begin{code}
    #include <stdio.h>

    int main() {
        int zahl = 42;
        printf("Die Zahl ist: %d\n", zahl);
        printf("Adresse: %p\n", (void*)&zahl);
        return 0;
    }
    \end{code}

    \textbf{Rückgabewert:}
    Die Anzahl der ausgegebenen Zeichen, oder \codeinline{-1} bei Fehler.

    \textbf{Siehe auch:} \codeinline{scanf()}, \codeinline{sprintf()}, \codeinline{fprintf()}
\end{man}



\subsection{C++}
\label{sec:reference:cpp}

\subsection{CMake}
\label{sec:reference:cmake}



\bibliographystyle{unsrt}
\bibliography{c-kurs}

\end{document}
